\chapter{Phosphate (Blood)}

\section*{What Is This Test?}
The phosphate blood test measures the concentration of phosphate (PO\textsubscript{4}\textsuperscript{3-}) in the blood serum or plasma. In clinical practice, ``serum phosphorus'' and ``serum phosphate'' are used interchangeably, though the test measures phosphorus in the form of phosphate. Phosphorus is the second most abundant mineral in the human body after calcium, with approximately 85\% of total body phosphorus (700--800 grams in adults) stored in bone and teeth as hydroxyapatite [Ca\textsubscript{10}(PO\textsubscript{4})\textsubscript{6}(OH)\textsubscript{2}]. The remaining 15\% is distributed in soft tissues (10\%), and less than 1\% in the extracellular fluid, making serum phosphate a partial but not complete indicator of total body stores. In blood, phosphate exists primarily as H\textsubscript{2}PO\textsubscript{4}\textsuperscript{-} and HPO\textsubscript{4}\textsuperscript{2-} (in a ratio of approximately 1:4 at pH 7.4), with approximately 10--15\% protein-bound. Phosphate is critical for: (1) bone mineralization as hydroxyapatite, (2) energy metabolism as the high-energy phosphate bonds in ATP, ADP, and creatine phosphate, (3) nucleic acid structure (DNA and RNA sugar-phosphate backbone), (4) cell membrane structure as phospholipids, (5) intracellular signaling as second messenger (cAMP, IP3), (6) regulation of hemoglobin oxygen affinity via 2,3-bisphosphoglycerate (2,3-BPG) in red blood cells, (7) acid-base buffering in urine (titratable acid as H\textsubscript{2}PO\textsubscript{4}\textsuperscript{-}), and (8) enzyme regulation by phosphorylation/dephosphorylation. Phosphate homeostasis is regulated by parathyroid hormone (PTH), fibroblast growth factor 23 (FGF23), and 1,25-dihydroxyvitamin D (calcitriol) \cite{berndt2007phosphatonins}. PTH decreases serum phosphate by inhibiting renal proximal tubular phosphate reabsorption (via sodium-phosphate cotransporter NaPi-IIa internalization) and stimulates bone resorption (which releases phosphate). FGF23 (produced by osteocytes) is the major phosphaturic hormone: it decreases renal phosphate reabsorption and suppresses calcitriol synthesis. Calcitriol increases both intestinal phosphate absorption and renal phosphate reabsorption. Phosphate has a significant diurnal variation: levels are highest in the early morning (nadir around 8--11 AM) and lowest in the afternoon, with a normal fluctuation of 0.5--1.0 mg/dL. Levels also decrease after meals due to insulin-mediated cellular phosphate uptake.

\section*{Alternative Names}
Phosphorus; Serum Phosphorus; Inorganic Phosphate; PO4; Serum Phosphate; Pi; P; Blood Phosphate Level

\section*{Principle}
Serum phosphate is measured by spectrophotometric methods based on the reaction of phosphate with ammonium molybdate to form phosphomolybdate complex. The most common method: phosphate + ammonium molybdate in acidic solution $\rightarrow$ ammonium phosphomolybdate (colorless) $\rightarrow$ reduced by ascorbic acid or aminonaphtholsulfonic acid (Fiske-SubbaRow reducing agent) $\rightarrow$ molybdenum blue (intense blue chromophore absorbing at 340 nm or 660--700 nm, proportional to phosphate concentration). This is the Fiske-SubbaRow method, first described in 1925 and still the basis of most modern automated phosphate assays \cite{tietz2012textbook}. Alternative reduction methods use stannous chloride or ferrous ion. Direct UV measurement of the unreduced phosphomolybdate at 340 nm is used on some analyzers without reduction to molybdenum blue. Enzymatic methods using phosphorylase, glucose oxidase, and peroxidase coupled reactions are used on some dry-chemistry platforms (e.g., Vitros). The test requires careful sample handling because serum phosphate rises over time as organic phosphate esters in red blood cells are hydrolyzed, releasing inorganic phosphate. This test is performed on serum separator tubes using spectrophotometric methods, not by hematology analyzers.

\section*{History}
Phosphorus was discovered by Hennig Brand in 1669 while attempting to create the philosopher's stone from urine. The essential role of phosphate in rickets was recognized by the early 20th century. The Fiske-SubbaRow method for phosphate measurement was developed in 1925 by Cyrus Fiske and Yellapragada SubbaRow at Harvard Medical School and remains the basis of modern methods. The discovery of PTH by James Collip in 1925 and its role in phosphate metabolism was subsequently established. FGF23 was discovered in 2000 through studies of autosomal dominant hypophosphatemic rickets (ADHR), revolutionizing the understanding of phosphate homeostasis. The identification of sodium-phosphate cotransporters (NaPi-IIa, NaPi-IIc) and their regulation by PTH and FGF23 occurred in the 1990s--2000s.

\section*{Why This Test Is Ordered}
\begin{itemize}
  \item Evaluation of symptoms consistent with hypophosphatemia: muscle weakness (including respiratory muscle weakness and difficulty weaning from mechanical ventilation), bone pain, osteomalacia, rhabdomyolysis, hemolytic anemia (ATP depletion impairs red cell membrane integrity and glycolysis), platelet dysfunction and bleeding tendency, metabolic encephalopathy (confusion, seizures, coma), cardiac dysfunction (decreased myocardial contractility)
  \item Evaluation of symptoms consistent with hyperphosphatemia: symptoms are primarily those of the resulting hypocalcemia---tetany, muscle cramps, paresthesias, seizures---or metastatic calcification in soft tissues, blood vessels, and joints (calciphylaxis in ESRD patients)
  \item Diagnosis and monitoring of parathyroid disorders: primary hyperparathyroidism causes hypophosphatemia (low phosphate + high calcium); hypoparathyroidism causes hyperphosphatemia (high phosphate + low calcium)
  \item Assessment of chronic kidney disease: hyperphosphatemia develops when GFR falls below 20--30 mL/min. Monitoring phosphate is essential because hyperphosphatemia drives secondary hyperparathyroidism, renal osteodystrophy, and vascular calcification. KDIGO guidelines recommend maintaining phosphate in the normal range in CKD stages 3--5D
  \item Monitoring of tumor lysis syndrome: massive phosphate release from lysed malignant cells (especially in Burkitt lymphoma, ALL, and other rapidly proliferating hematologic malignancies after chemotherapy). Hyperphosphatemia can cause acute kidney injury from calcium-phosphate crystal precipitation in renal tubules
  \item Evaluation of rhabdomyolysis: muscle cell destruction releases phosphate, and acute kidney injury impairs excretion, causing severe hyperphosphatemia
  \item Monitoring refeeding syndrome: after initiating nutrition in severely malnourished patients (anorexia nervosa, chronic alcoholism, prolonged starvation), insulin-mediated cellular uptake of phosphate, potassium, and magnesium can cause severe hypophosphatemia (<1.0 mg/dL) within 24--72 hours of refeeding. The hallmark triad is hypophosphatemia, hypokalemia, and hypomagnesemia
  \item Evaluation of suspected oncogenic osteomalacia (tumor-induced osteomalacia): rare mesenchymal tumors secreting FGF23 cause severe hypophosphatemia, renal phosphate wasting, low 1,25(OH)\textsubscript{2}D, and osteomalacia
  \item Monitoring of parenteral nutrition (TPN): patients receiving TPN without adequate phosphate supplementation rapidly develop hypophosphatemia
  \item Diabetic ketoacidosis (DKA): total body phosphate is depleted, though serum phosphate may be normal or elevated initially due to transcellular shifts in acidosis. With insulin therapy, phosphate shifts intracellularly, and severe hypophosphatemia can develop during DKA treatment
\end{itemize}

\section*{Primary vs Secondary Test}
Serum phosphate is a primary screening test as part of the comprehensive metabolic panel (CMP) in many laboratories, though some CMP configurations omit phosphate. It is routinely measured in patients with CKD, suspected parathyroid disorders, malnutrition, and critical illness \cite{tierney2023current}.

\section*{How This Test Is Performed}
A healthcare professional draws blood from a vein into a serum separator tube (gold-top) or lithium heparin tube (green-top). The patient should be fasting (8 hours) for optimal accuracy, as serum phosphate decreases postprandially due to insulin-mediated cellular uptake. The sample must be centrifuged and separated from cells within 1--2 hours of collection. If serum remains in contact with red blood cells for prolonged periods, organic phosphate esters are hydrolyzed by cellular enzymes, releasing inorganic phosphate and artifactually elevating measured levels. The sample is analyzed on automated chemistry analyzers using the phosphomolybdate method \cite{fischbach2023manual}. Results available within 1--2 hours.

\section*{What Preparation Is Needed}
Fasting for 8 hours is strongly recommended because food intake decreases serum phosphate by 0.5--1.0 mg/dL due to insulin-mediated cellular phosphate uptake. The sample should be drawn in the morning due to diurnal variation (levels highest in early morning). Inform the provider of medications affecting phosphate: phosphate-binding antacids (calcium carbonate, aluminum hydroxide, sevelamer, lanthanum---decrease phosphate absorption), IV glucose/insulin therapy (shifts phosphate intracellularly), oral phosphate supplements, calcitriol or vitamin D therapy, bisphosphonates, and diuretics (acetazolamide increases urinary phosphate). Samples must be processed promptly---delays >2 hours before serum separation cause significant artifactual elevation from red cell phosphate ester hydrolysis. Hemolyzed samples must be rejected (red cell phosphate is 6--8 times higher than serum phosphate).

\section*{Contraindications and Precautions}
\begin{itemize}
  \item No absolute contraindications. Critical pre-analytical considerations:
  \item (1) Hemolysis: the most critical interferent---intracellular phosphate concentration is 30--40 mg/dL (6--8$\times$ serum); even microhemolysis can significantly elevate measured phosphate. Hemolyzed samples must be rejected.
  \item (2) Delayed serum separation: phosphate esters in red blood cells are hydrolyzed by phosphatases over time, releasing inorganic phosphate that falsely elevates the measured value by 1--2 mg/dL after 4--6 hours.
  \item (3) Postprandial state: eating decreases serum phosphate by 0.5--1.0 mg/dL; fasting samples are essential for accurate baseline.
  \item (4) Diurnal variation: phosphate is 0.5--1.0 mg/dL lower in the afternoon than in the morning.
  \item (5) EDTA tubes: EDTA interferes with the phosphomolybdate reaction and EDTA chelates magnesium, which is a cofactor in some enzymatic methods. Use serum separator tubes.
  \item (6) Medications affecting results: phosphate binders (decrease), IV glucose/insulin (rapid decrease), mannitol (interferes with some assays), liposomal amphotericin B formulations (contain phospholipid that is measured as phosphate).
  \item (7) Hyperlipidemia: severe lipemia can cause turbidity that interferes with spectrophotometric methods.
  \item (8) Contrast dye and certain drugs can interfere colorimetrically.
  \item (9) Pregnancy: serum phosphate is slightly lower due to increased renal phosphate clearance and fetal demands.
\end{itemize}
\section*{Risks}
Minimal risk from venipuncture. Mild pain or bruising at the puncture site resolves in 1--2 days. Rare: hematoma, infection, nerve injury, vasovagal syncope. Blood volume drawn (3--5 mL) is negligible.

\section*{What Happens After the Test}
Results available in 1--2 hours. Critical values: phosphate <1.0 mg/dL (risk of respiratory failure, rhabdomyolysis, hemolytic anemia, cardiac dysfunction) or >9.0 mg/dL (risk of acute kidney injury from calcium-phosphate precipitation, fatal hypocalcemia) \cite{wallach2022interpretation}. Phosphate <1.0 mg/dL: severe hypophosphatemia requiring IV phosphate replacement (sodium phosphate or potassium phosphate at 2.5--5.0 mg/kg over 6 hours, monitoring calcium $\times$ phosphate product to avoid precipitation). Phosphate >5.5 mg/dL in CKD patients: manage with dietary phosphate restriction (800--1000 mg/day), phosphate binders (calcium acetate, sevelamer, lanthanum, ferric citrate, sucroferric oxyhydroxide), and adequate dialysis \cite{shaman2016hyperphosphatemia}. The calcium $\times$ phosphate product (Ca $\times$ PO4 in mg/dL) should be kept <55 mg\textsuperscript{2}/dL\textsuperscript{2} in CKD patients to minimize vascular calcification risk.

\section*{Understanding Results}

\subsection*{Below Normal}
\begin{itemize}
  \item \textbf{Hypophosphatemia (PO4 <2.5 mg/dL or <0.81 mmol/L):} Present in approximately 2--5\% of hospitalized patients; up to 20--30\% of ICU patients, especially those with sepsis, DKA, alcoholism, or receiving mechanical ventilation. Mild 1.5--2.5 mg/dL; Moderate 1.0--1.5 mg/dL; Severe <1.0 mg/dL \cite{amanzadeh2006hypophosphatemia}.
  \item \textbf{Renal phosphate wasting (most common chronic cause):} Primary hyperparathyroidism (elevated PTH downregulates NaPi-IIa, causing renal phosphate wasting---labs show high Ca, low PO4, high/normal PTH, normal/elevated calcitriol). Vitamin D deficiency (rickets/osteomalacia---low calcium and phosphate absorption, secondary hyperparathyroidism adds renal wasting; labs show low Ca, low PO4, elevated PTH, low 25(OH)D). Oncogenic osteomalacia (FGF23-secreting mesenchymal tumors---severe phosphaturic hypophosphatemia, low 1,25(OH)\textsubscript{2}D, normal Ca, normal/low PTH). X-linked hypophosphatemic rickets (XLH---mutation in PHEX gene causing elevated FGF23). Autosomal dominant hypophosphatemic rickets (ADHR---mutation in FGF23 cleavage site). Fanconi syndrome (generalized proximal tubular dysfunction with phosphate, glucose, amino acid, uric acid, and bicarbonate wasting; causes: cystinosis, Wilson disease, heavy metals, tenofovir, outdated tetracycline, multiple myeloma, Sj\"{o}gren's syndrome). Renal tubular acidosis (especially type II proximal RTA). Post-renal transplant (persistent hyperparathyroidism and glucocorticoid-induced phosphaturia). Diuretics (acetazolamide).
  \item \textbf{Transcellular shifts (most common acute cause in hospitalized patients):} Refeeding syndrome (insulin-mediated cellular phosphate uptake within 24--72 hours of initiating nutrition in severely malnourished patients; define risk: BMI <16, unintentional weight loss >15\% in 3--6 months, little or no nutritional intake >10 days, low baseline K/Mg/PO4). Diabetic ketoacidosis treatment (insulin drives phosphate into cells; despite total body depletion, serum PO4 may be normal at presentation due to transcellular shifts from acidosis; falls precipitously with insulin therapy---monitor every 4--6 hours during DKA treatment). Respiratory alkalosis (hyperventilation, mechanical ventilation, salicylate poisoning, sepsis, anxiety---increased intracellular pH activates phosphofructokinase, accelerating glycolysis and cellular phosphate consumption; most common cause of hypophosphatemia in mechanically ventilated patients). IV glucose/insulin administration (standard DKA or hyperkalemia treatment). Hungry bone syndrome (calcium and phosphate deposition into bone after parathyroidectomy for severe hyperparathyroidism). Catecholamine infusions (beta-adrenergic stimulation shifts phosphate intracellularly). Therapeutic hypothermia.
  \item \textbf{Decreased intestinal absorption:} Severe malnutrition, chronic alcoholism (poor intake + magnesium deficiency impairing renal phosphate conservation + respiratory alkalosis from alcohol withdrawal), malabsorption (celiac disease, Crohn's, short bowel), phosphate binders (calcium carbonate, aluminum hydroxide, sevelamer, lanthanum), vitamin D deficiency, chronic antacid use.
  \item \textbf{Clinical manifestations by severity:} Moderate (1.0--2.5 mg/dL): fatigue, bone pain, muscle weakness. Severe (<1.0--1.5 mg/dL): muscle weakness (including diaphragmatic weakness and failure to wean from ventilator), rhabdomyolysis, hemolytic anemia (ATP depletion reduces RBC deformability and Na\textsuperscript{+}/K\textsuperscript{+}-ATPase function $\rightarrow$ hemolysis), leukocyte dysfunction (impaired chemotaxis and phagocytosis due to ATP depletion), thrombocytopenia and platelet dysfunction, metabolic encephalopathy (irritability, confusion, seizures, coma), and impaired cardiac contractility (reversible cardiomyopathy). Very severe (<0.5 mg/dL): usually fatal without aggressive IV replacement.
\end{itemize}

\subsection*{Above Normal}
\begin{itemize}
  \item \textbf{Hyperphosphatemia (PO4 >4.5 mg/dL or >1.45 mmol/L in adults):} Primarily a problem of impaired renal excretion, since the kidney can normally excrete large phosphate loads. Mild 4.5--5.5 mg/dL; Moderate 5.5--7.0 mg/dL; Severe >7.0 mg/dL.
  \item \textbf{Renal failure (most common cause):} Acute kidney injury; chronic kidney disease (GFR <20--30 mL/min impairs phosphate excretion, leading to phosphate retention, secondary hyperparathyroidism from PTH stimulation, and reduced calcitriol synthesis from kidney damage and FGF23 elevation). Hyperphosphatemia drives progression of CKD and is independently associated with increased cardiovascular mortality in ESRD patients through vascular calcification \cite{kestenbaum2007phosphate}.
  \item \textbf{Hypoparathyroidism:} Post-surgical (thyroidectomy/parathyroidectomy), autoimmune, congenital (DiGeorge syndrome), activating CaSR mutation. Labs: low Ca, high PO4, low PTH. The elevated phosphate results from increased renal phosphate reabsorption in the absence of PTH.
  \item \textbf{Pseudohypoparathyroidism:} Target organ resistance to PTH. Labs: low Ca, high PO4, elevated PTH (distinguishing from hypoparathyroidism where PTH is low). Associated with Albright hereditary osteodystrophy (short stature, round face, short fourth metacarpals, subcutaneous ossifications) in type 1a (GNAS1 mutation).
  \item \textbf{Tumor lysis syndrome:} Massive cell lysis after chemotherapy for rapidly proliferating hematologic malignancies (Burkitt lymphoma, ALL, high-grade NHL) releases intracellular phosphate (along with potassium, uric acid, and LDH). Hyperphosphatemia causes secondary hypocalcemia from calcium-phosphate precipitation, which can lead to acute kidney injury from crystal nephropathy, tetany, seizures, and cardiac arrhythmias. Prophylaxis with aggressive IV hydration, rasburicase (for uric acid), and close monitoring of phosphate and calcium is standard.
  \item \textbf{Rhabdomyolysis:} Muscle cell destruction releases phosphate, potassium, and CK. If AKI develops, phosphate excretion is further impaired, resulting in severe hyperphosphatemia.
  \item \textbf{Excessive intake:} Oral phosphate supplements (sodium phosphate bowel preparations for colonoscopy---can cause acute phosphate nephropathy, especially in elderly, CKD, or volume-depleted patients), phosphate-containing enemas (especially Fleet enemas in children or patients with renal impairment), IV phosphate administration, vitamin D toxicity (increases intestinal phosphate absorption).
  \item \textbf{Transcellular shifts:} Acidosis (especially during DKA or lactic acidosis---intracellular phosphate shifts extracellularly), severe hemolysis, tumor necrosis, catabolic states, bisphosphonate therapy.
  \item \textbf{Other causes:} Acromegaly (GH increases renal phosphate reabsorption), hyperthyroidism, bisphosphonate therapy (transient hyperphosphatemia with etidronate due to increased renal reabsorption), magnesium deficiency (unclear mechanism, mild hyperphosphatemia).
  \item \textbf{Clinical manifestations:} Acute hyperphosphatemia causes hypocalcemia through calcium-phosphate precipitation (Ca $\times$ PO4 product >55--70 mg\textsuperscript{2}/dL\textsuperscript{2}). Symptoms are those of resultant hypocalcemia: muscle cramps, tetany, paresthesias, Chvostek and Trousseau signs, prolonged QT, seizures. Chronic hyperphosphatemia in CKD causes secondary hyperparathyroidism, renal osteodystrophy (high turnover bone disease, osteitis fibrosa cystica), vascular calcification (medial calcification of coronary and peripheral arteries), and calciphylaxis (calcific uremic arteriolopathy---painful ischemic skin necrosis with high mortality).
\end{itemize}

\section*{Normal Results}
\begin{itemize}
  \item \textbf{Adults:} 2.5--4.5 mg/dL (0.81--1.45 mmol/L)
  \item \textbf{Children (1--12 years):} 3.5--5.5 mg/dL (higher due to growth hormone effects and higher bone turnover)
  \item \textbf{Infants (1--12 months):} 4.5--6.5 mg/dL
  \item \textbf{Newborns (0--7 days):} 4.3--9.3 mg/dL (highest in first week of life)
  \item \textbf{CKD stage 3--5 target:} 2.5--4.5 mg/dL (KDIGO guidelines)
  \item \textbf{Critical low value:} <1.0 mg/dL
  \item \textbf{Critical high value:} >9.0 mg/dL
  \item \textbf{Calcium $\times$ phosphate product:} <55 mg\textsuperscript{2}/dL\textsuperscript{2} (target in CKD)
  \item \textbf{Conversion:} 1 mg/dL = 0.323 mmol/L; 1 mmol/L = 3.1 mg/dL
\end{itemize}

\section*{Follow-up Tests}
\begin{itemize}
  \item \textbf{Calcium (total and ionized):} Essential to assess concurrent calcium disturbance. Hyperphosphatemia drives hypocalcemia through precipitation and PTH-vitamin D axis disruption.
  \item \textbf{Parathyroid hormone (PTH, intact):} Differentiates PTH-mediated phosphate disorders. High Ca + low PO4 + elevated PTH $\rightarrow$ primary hyperparathyroidism. Low Ca + high PO4 + low PTH $\rightarrow$ hypoparathyroidism. Low Ca + high PO4 + elevated PTH $\rightarrow$ pseudohypoparathyroidism or vitamin D deficiency.
  \item \textbf{25-hydroxyvitamin D and 1,25-dihydroxyvitamin D:} Vitamin D deficiency causes hypocalcemia and hypophosphatemia with secondary hyperparathyroidism. Low 1,25(OH)\textsubscript{2}D with hypophosphatemia and normal 25(OH)D suggests FGF23-mediated phosphate wasting (oncogenic osteomalacia, XLH).
  \item \textbf{24-hour urine phosphate and fractional excretion of phosphate (FEPO4):} In hypophosphatemia, FEPO4 >5\% (or 24h urine PO4 >100 mg) indicates inappropriate renal phosphate wasting. FEPO4 <5\% in hypophosphatemia indicates appropriate renal conservation (GI loss, refeeding, transcellular shift).
  \item \textbf{FGF23 (fibroblast growth factor 23):} Order if oncogenic osteomalacia or X-linked hypophosphatemia is suspected. Elevated FGF23 causes hypophosphatemia, phosphaturia, and low calcitriol. Also elevated in CKD as a compensatory response to phosphate retention.
  \item \textbf{Renal function (BUN, creatinine, eGFR):} Essential because hyperphosphatemia is almost always due to impaired renal excretion.
  \item \textbf{Magnesium level:} Hypomagnesemia can cause mild hypophosphatemia and frequently accompanies hypophosphatemia in refeeding syndrome, alcoholism, and DKA.
  \item \textbf{Creatine kinase (CK):} For suspected rhabdomyolysis causing hyperphosphatemia.
  \item \textbf{Uric acid and LDH:} For suspected tumor lysis syndrome.
  \item \textbf{ECG:} For QT changes from associated hypocalcemia or direct effects of severe phosphate disturbance.
  \item \textbf{Bone-specific alkaline phosphatase (BSAP):} For assessment of bone turnover in renal osteodystrophy.
  \item \textbf{Renal ultrasound:} May detect nephrocalcinosis from chronic hyperphosphatemia or nephrolithiasis from chronic hypercalciuria.
\end{itemize}

\section*{References}
\bibliographystyle{unsrtnat}
\bibliography{references}

\clearpage
\section*{Regulatory Status and Authorization Date}
This chapter describes a test category, not a specific commercial product. FDA, CDSCO, EU IVDR, and MHRA approval or registration dates therefore cannot be assigned without the assay manufacturer, product name, instrument, and intended use. For laboratory-developed tests, an individual agency approval date may not exist. See the Preface for the interpretation of regulatory status.

\clearpage
